\chapter{How Bonds Trade}\label{ch:m2:how-bonds-trade}

A large company may have fifty bonds outstanding and its stock one line on an
exchange. Each bond is a separate security with its own coupon and maturity, and
most of them do not trade on most days. There is no order book to consult; an
investor who wants to sell asks dealers for prices, and the answer depends on who
is asked, how many, and what they learn from being asked. Until 2002 even the
prices of past trades were private. Since then, American dealers have reported
every corporate bond trade to a public tape, and since 15 May 2023 they have
flagged trades done as part of a portfolio: in its first week, those were already
5.2\% of the notional traded. This chapter explains the \omterm{def:m2:how-bonds-trade:rfq}{request for quote}, \omterm{def:m2:how-bonds-trade:trace}{trade
reporting}, the new ways bonds trade in bulk or with each other, the bond
exchange-traded funds that sit on top of the market, and how dealers price what
rarely trades.

\section{Request for quote}

\begin{definition}[Request for quote, axe]\label{def:m2:how-bonds-trade:rfq}
A \emph{request for quote}\index{request for quote} (RFQ) is a client's request,
usually sent electronically to several dealers at once, for a firm price to buy or
sell a stated amount of a security; the client trades with the best answer or
none. An \emph{axe}\index{axe} is a dealer's advertised interest in buying or
selling a particular bond, usually because of its inventory.
\end{definition}

The RFQ is an auction the client runs for each trade. Asking more dealers brings
more competition and a better best price. But each dealer asked learns that someone
wants to sell that bond now, and a losing dealer may use that knowledge, selling
the bond or its hedges ahead of the winner, or simply marking it down; the client
pays for that through the price of this trade or the next. \omterm{def:m2:how-bonds-trade:rfq}{Axes} tell the client
which dealers want the other side and so are likely to bid well.

\begin{proposition}[How many dealers to ask]\label{prop:m2:how-bonds-trade:n}
Suppose dealer $i$ bids $v - m + \sigma Z_i$, with $v$ the bond's value, $m$ an average
markup and $Z_i$ independent standard normal private terms (inventory, \omterm{def:m2:how-bonds-trade:rfq}{axes}, views),
and that each losing dealer costs the client $\ell$ through leakage. Selling to the
best of $n$ bids costs, in expectation,
\[
C(n) \;=\; m - \sigma\,e_n + \ell\,(n-1),
\]
where $e_n$ is the expected maximum of $n$ standard normals ($e_1 = 0$, $e_2 = 0.564$,
$e_3 = 0.846$, $e_4 = 1.029$). Since $e_n$ grows ever more slowly, $C$ has a
minimum at a finite $n$, which falls as the leakage $\ell$ rises.
\end{proposition}

\begin{proof}
The best bid is $v - m + \sigma\max_i Z_i$, whose expected discount to $v$ is
$m - \sigma e_n$; the $n-1$ losers add $\ell$ each. The increments $e_{n+1} - e_n$
decrease to zero, while leakage adds $\ell$ for every dealer, so $C(n+1) - C(n) =
\ell - \sigma(e_{n+1} - e_n)$ turns positive.
\end{proof}

\begin{omfigure}
\begin{tikzpicture}
\begin{axis}[width=11cm,height=4.8cm,
  xlabel={dealers asked, $n$},ylabel={expected cost (bp)},
  tick label style={font=\small},label style={font=\small},
  xmin=0.7,xmax=8.3,ymin=5.5,ymax=14,grid=major,grid style={gray!25},
  legend columns=3,legend style={font=\footnotesize,at={(0.5,-0.3)},anchor=north,draw=none,fill=none,/tikz/every even column/.append style={column sep=8pt}},
  table/col sep=comma]
\addplot[omDef,very thick,mark=*,mark size=1.4pt] table[x=n,y=l05]{figdata/markets-2/22-how-bonds-trade/cost.csv};
\addplot[omThm,very thick,mark=*,mark size=1.4pt] table[x=n,y=l10]{figdata/markets-2/22-how-bonds-trade/cost.csv};
\addplot[omProp,very thick,mark=*,mark size=1.4pt] table[x=n,y=l15]{figdata/markets-2/22-how-bonds-trade/cost.csv};
\legend{leakage 0.5 bp,leakage 1.0 bp,leakage 1.5 bp}
\end{axis}
\end{tikzpicture}

\omcaption{A client's expected cost of selling a bond by RFQ, in basis points of
price, by the number of dealers asked, with an average markup of 10 basis points, a
spread of private terms of 5 and three levels of leakage per losing dealer. The best
number of dealers is 6, 3 and 2. Illustrative; data: the chapter's
tutorial.}\label{fig:m2:how-bonds-trade:cost}
\end{omfigure}

\begin{example}[Three dealers]\label{ex:m2:how-bonds-trade:three}
With a markup of 10 basis points, private terms spread by 5 and a leakage of 1 basis
point per loser, asking one dealer costs 10.0 basis points, two 8.18, three 7.77,
four 7.85 and six 8.66. Three is best; on a USD~25 million sale it costs about
USD~19\,400, against USD~25\,000 with a single dealer.
\end{example}

\begin{omfigure}
\begin{tikzpicture}
\begin{axis}[width=11cm,height=4.8cm,
  xlabel={winning bid's discount to value (bp)},ylabel={density},
  tick label style={font=\small},label style={font=\small},
  xmin=-10,xmax=26,ymin=0,ymax=0.14,grid=major,grid style={gray!25},
  yticklabel style={/pgf/number format/.cd,fixed,precision=2},
  legend columns=3,legend style={font=\footnotesize,at={(0.5,-0.3)},anchor=north,draw=none,fill=none,/tikz/every even column/.append style={column sep=8pt}},
  table/col sep=comma]
\addplot[omDef,thick,const plot mark mid] table[x=x,y=n1]{figdata/markets-2/22-how-bonds-trade/hist.csv};
\addplot[omThm,very thick,const plot mark mid] table[x=x,y=n3]{figdata/markets-2/22-how-bonds-trade/hist.csv};
\addplot[omProp,thick,dashed,const plot mark mid] table[x=x,y=n6]{figdata/markets-2/22-how-bonds-trade/hist.csv};
\legend{1 dealer,3 dealers,6 dealers}
\end{axis}
\end{tikzpicture}

\omcaption{Simulated distribution of the winning bid's discount to the bond's value,
before leakage, when one, three or six dealers are asked. More dealers shift the
distribution towards better prices and narrow it, with diminishing gains.
Illustrative; data: the chapter's tutorial.}\label{fig:m2:how-bonds-trade:hist}
\end{omfigure}

\section{Trade reporting and transparency}

\begin{definition}[Trade reporting, dissemination cap]\label{def:m2:how-bonds-trade:trace}
\emph{Trade reporting}\index{trade reporting} is the obligation of a dealer to report
each trade, with its price, size and time, to a regulator's system that publishes
it. A \emph{dissemination cap}\index{dissemination cap} is a size above which the
published report shows only that the trade was larger than the cap.
\end{definition}

FINRA's TRACE system has collected American corporate bond trades since July 2002;
dealers must report within fifteen minutes, and the reports are published as they
arrive. Publication came in phases over three years, the largest, most highly rated
bonds first, and large trades have been published with capped sizes: exact amounts
up to USD~5 million for investment grade and USD~1 million for high yield, ``5MM$+$''
and ``1MM$+$'' above. The cap protects a dealer who has bought a large
block from being front-run while it resells it; the published price still tells
the market where bonds trade. The European Union had no consolidated tape of its bond
trades; in July 2025 its securities regulator selected the first provider of one.

\begin{dated}{2026-09}{Transparency rules}\label{dat:m2:how-bonds-trade:transparency}
United States: TRACE reporting within 15 minutes for corporate bonds;
portfolio-trade modifier since 15 May 2023. European Union: ESMA selected Ediphy
(fairCT) on 3 July 2025 as the first consolidated tape provider for bonds, to operate
for five years after authorisation.
\end{dated}

\section{All-to-all and portfolio trading}

\begin{definition}[All-to-all trading, portfolio trade]\label{def:m2:how-bonds-trade:a2a}
\emph{All-to-all trading}\index{all-to-all trading} lets any participant on a
platform, investor or dealer, answer any other's request or order, instead of
investors asking only dealers. A \emph{portfolio trade}\index{portfolio trade} is the
purchase or sale of a basket of bonds at a single agreed price for the whole basket;
TRACE flags a trade as such when it is between two parties, for at least ten
different bonds at one price.
\end{definition}

\omterm{def:m2:how-bonds-trade:a2a}{All-to-all trading} widens the set of possible buyers: an asset manager selling a
bond may find another manager who wants it, with no dealer in between. Portfolio
trading does the opposite of RFQ for single bonds: instead of asking for a hundred
prices, the client asks for one, for the whole list, and the dealer prices the risk of
the basket, which it can hedge with credit indices and bond ETFs rather than bond by
bond. The price of each bond in a \omterm{def:m2:how-bonds-trade:a2a}{portfolio trade} is not a negotiated price of that
bond, which is why FINRA required it to be flagged: a reader of the tape should not
take it as one.

\begin{omfigure}
\centering
\begin{tikzpicture}[font=\small,
  box/.style={draw,thick,rounded corners=2pt,align=center,minimum height=1cm,minimum width=2.6cm,inner sep=3pt},
  arr/.style={-{Stealth[length=2.5mm]},thick}]
\node[box,draw=omDef,fill=omDef!8] (cl) at (0,0) {client with a\\list of 150 bonds};
\node[box,draw=omThm,fill=omThm!8] (dl) at (5.5,0) {portfolio-trading\\dealer};
\node[box,draw=omProp,fill=omProp!8] (etf) at (11,1.1) {bond ETF\\(creations, redemptions)};
\node[box,draw=gray,fill=gray!8] (cdx) at (11,-1.1) {credit index\\swaps};
\draw[arr] (cl) -- node[above,font=\footnotesize,align=center] {one price\\for the list} (dl);
\draw[arr] (dl) -- node[above,sloped,font=\footnotesize] {hedge} (etf);
\draw[arr] (dl) -- node[below,sloped,font=\footnotesize] {hedge} (cdx);
\node[font=\footnotesize,text=omThm,align=center] at (5.5,-1.4) {works the bonds out\\over days};
\end{tikzpicture}

\omcaption{A \omterm{def:m2:how-bonds-trade:a2a}{portfolio trade}. The client asks one price for a whole list; the dealer
takes the basket, hedges its overall risk with bond ETFs and index swaps, and sells or
redeems the individual bonds over the following days. Schematic.}\label{fig:m2:how-bonds-trade:pt}
\end{omfigure}

\section{Bond ETFs as a liquidity layer}

A bond exchange-traded fund (One Quant Book 1, chapter 14) holds hundreds or thousands
of bonds and trades on an exchange all day. Its shares trade far more easily than
most of its bonds, and its authorised participants create and redeem shares against
baskets of bonds: when the ETF trades below the value of its bonds, they buy shares,
redeem them for bonds and sell the bonds; when it trades above, they do the opposite.
The ETF is thus a liquid layer on top of an illiquid market: an investor can change
its credit exposure in minutes, and the heavy lifting of trading the bonds is left to
the participants, in baskets (\cref{fig:m2:how-bonds-trade:etf}). In stressed markets
the ETF's price can fall below the published value of its bonds, whose prices are
stale because they have not traded, and the ETF becomes the place where prices are
discovered.

\begin{omfigure}
\centering
\begin{tikzpicture}[font=\small,
  box/.style={draw,thick,rounded corners=2pt,align=center,minimum height=1cm,minimum width=2.6cm,inner sep=3pt},
  arr/.style={{Stealth[length=2.5mm]}-{Stealth[length=2.5mm]},thick}]
\node[box,draw=omDef,fill=omDef!8] (inv) at (0,0) {investors\\(trade shares\\all day)};
\node[box,draw=omProp,fill=omProp!8] (ex) at (4.3,0) {exchange:\\ETF shares};
\node[box,draw=omThm,fill=omThm!8] (ap) at (8.6,0) {authorised\\participants};
\node[box,draw=gray,fill=gray!8] (fund) at (8.6,-2.6) {the fund\\(creations, redemptions)};
\node[box,draw=gray,fill=gray!8] (otc) at (12.9,0) {bond market\\(RFQ, portfolios)};
\draw[arr] (inv) -- (ex); \draw[arr] (ex) -- (ap); \draw[arr] (ap) -- node[right,font=\footnotesize,align=left] {shares for\\baskets} (fund);
\draw[arr] (ap) -- node[above,font=\footnotesize] {bonds} (otc);
\end{tikzpicture}

\omcaption{A bond ETF as a liquidity layer. Investors trade shares on an exchange;
authorised participants exchange shares for baskets of bonds with the fund, and buy
or sell the bonds in the over-the-counter market, keeping the share price near the
value of the bonds. Schematic.}\label{fig:m2:how-bonds-trade:etf}
\end{omfigure}

\section{Algorithmic dealer pricing}

\begin{definition}[Composite price]\label{def:m2:how-bonds-trade:composite}
A \emph{composite price}\index{composite price} of a bond is a single estimated price
built from many observations, recent trades, dealer quotes and prices of related
bonds, weighted by their recency and reliability, with outliers removed.
\end{definition}

A dealer that answers thousands of RFQs a day cannot price each by hand. Its pricing
engine starts from a composite for each bond, anchored on the last trades reported to
TRACE and on quotes, moves it with the issuer's other bonds, its CDS and the
relevant index, and then adds a markup that depends on the size, the direction, the
client, and the dealer's own inventory and \omterm{def:m2:how-bonds-trade:rfq}{axes}. The same engine sets how many RFQs
to answer and how aggressively, the other side of
\cref{prop:m2:how-bonds-trade:n}.

\begin{example}[A composite]\label{ex:m2:how-bonds-trade:composite}
Quotes of 99.42, 99.45, 99.40 and 99.47, 10, 40, 90 and 5 seconds old, and a stale
print at 98.20, 20 seconds old: the outlier is dropped by the median-deviation rule,
and weighting the rest by a one-minute half-life gives a composite of 99.441.
\end{example}

\section{Tutorial: an RFQ simulator}\label{tut:m2:how-bonds-trade:rfq}

\begin{tutorial}
\textbf{Goal.} Simulate RFQ auctions, compute the client's expected cost by the number
of dealers with leakage, find the best number, and build a \omterm{def:m2:how-bonds-trade:composite}{composite price}.
\textbf{End state:} \cref{fig:m2:how-bonds-trade:cost,fig:m2:how-bonds-trade:hist},
\cref{ex:m2:how-bonds-trade:three,ex:m2:how-bonds-trade:composite} and the numbers of
the weekend problem.
\begin{tutsteps}
\item \textbf{The auction}: expected best bid, cost with leakage, the best $n$.
\omcode{code/firm/rfq/firm_rfq.py}{14}{40}{Expected best bid, client cost with leakage, best number of dealers, simulation.}\label{lst:m2:how-bonds-trade:rfq}
\item \textbf{The composite}: outliers removed, recency weights.
\omcode{code/firm/rfq/firm_rfq.py}{43}{51}{A composite price from quotes of different ages.}\label{lst:m2:how-bonds-trade:composite}
\item \textbf{Run} \texttt{rfq\_demo.problem()}, \texttt{rfq\_demo.composite\_example()}
and \texttt{fig\_rfq.py}.
\end{tutsteps}
\textbf{What to change next.} Make the leakage grow with the size of the trade and see
the best number of dealers fall for large trades; then give one dealer an \omterm{def:m2:how-bonds-trade:rfq}{axe} (a
higher mean bid) and find when asking only that dealer is best.
\end{tutorial}

\section{Build: the RFQ simulator}\label{bld:m2:how-bonds-trade:rfq}

\begin{build}
\bfield{Purpose} The miniature firm both asks for and answers RFQs in bonds: as a
client it chooses how many dealers to ask; as a dealer it prices from composites.
\bfield{Interface} \texttt{expected\_max\_normal(n)}; \texttt{expected\_cost(n, markup,
sigma, leak)}; \texttt{best\_n(markup, sigma, leak, n\_max)};
\texttt{simulate\_rfq(n, markup, sigma, trials, seed)}; \texttt{composite(quotes,
half\_life, k)}.
\bfield{Rules} Costs in basis points of price; independent normal private terms;
leakage linear in the number of losers; composite with median-deviation outlier
removal and exponential recency weights.
\bfield{Acceptance tests} \texttt{code/firm/rfq/tests/}: $e_1 = 0$ and $e_2 \approx
1/\sqrt\pi$; no leakage favours asking many, heavy leakage one; the simulated mean
matches the formula; the composite ignores an outlier.
\bfield{Stretch} Dealers who learn the client's history; correlated private terms;
a composite fitted across the issuer's curve; RFQ response models trained on
platform data (One Quant Book 12).
\end{build}

\begin{omsources}
\item FINRA, TRACE overview; Regulatory Notice 22-12 on portfolio trades.
\item P.~Asquith, T.~Covert and P.~Pathak, ``The effects of mandatory transparency in
financial market design: evidence from the corporate bond market'', NBER Working
Paper 19417, online appendix.
\item ICE, ``Early observations about the new portfolio trade flag from FINRA TRACE'',
May 2023.
\item ESMA, selection of the consolidated tape provider for bonds, July 2025.
\end{omsources}

\section{Exercises}

\begin{exercise}[$\star$]\label{exo:m2:how-bonds-trade:1}
A TRACE print shows an investment-grade bond traded ``5MM$+$'' at 99.50. What do you
know about the trade's size?
\end{exercise}

\begin{exercise}[$\star$]\label{exo:m2:how-bonds-trade:2}
With the chapter's parameters and no leakage, what is the expected cost of asking one,
two and four dealers?
\end{exercise}

\begin{exercise}[$\star$]\label{exo:m2:how-bonds-trade:3}
Why must \omterm{def:m2:how-bonds-trade:a2a}{portfolio trades} be flagged on the tape?
\end{exercise}

\begin{exercise}[$\star\star$]\label{exo:m2:how-bonds-trade:4}
Give the best number of dealers to ask when leakage is 0.5, 1.0 and 1.5 basis points
per loser.
\end{exercise}

\begin{exercise}[$\star\star$]\label{exo:m2:how-bonds-trade:5}
A bond ETF trades at 2\% below the value of its bonds in a stressed week. What does an
authorised participant do, and why might it hesitate?
\end{exercise}

\begin{exercise}[$\star\star$]\label{exo:m2:how-bonds-trade:6}
Why might a dealer decline to answer an RFQ at all?
\end{exercise}

\begin{exercise}[$\star\star\star$]\label{exo:m2:how-bonds-trade:7}
\emph{Coding.} With \texttt{composite}, recompute \cref{ex:m2:how-bonds-trade:composite}
with a half-life of 10 seconds, and explain the change.
\end{exercise}

\begin{exercise}[$\star\star\star$]\label{exo:m2:how-bonds-trade:8}
\emph{Find the flaw.} ``Always send the RFQ to every dealer on the platform: more
competition can only give a better price.'' Correct it.
\end{exercise}

\section{Problem: How Many Dealers?}

\begin{problem}[{Weekend problem --- the cost of competition and of being seen}]\label{pb:m2:how-bonds-trade:1}
An asset manager must sell USD~25 million of a corporate bond. Dealers' bids average 10
basis points below value, with private terms spread by 5 basis points; each dealer that
sees the request and loses costs the manager 1 basis point on average.

\medskip\noindent\textbf{Part I --- The auction.}
\begin{enumerate}
\item Give the expected cost of asking one dealer, in basis points and dollars.
\item Give the expected costs of two, three, four and six dealers.
\item What is the best number of dealers, and what does it cost?
\item What would it save against asking six?
\item Why does the gain from each extra dealer fall?
\end{enumerate}

\medskip\noindent\textbf{Part II --- Leakage.}
\begin{enumerate}[resume]
\item How does a losing dealer's knowledge cost the manager?
\item Give the best number of dealers if leakage is 0.5 and 1.5 basis points.
\item Why does leakage grow with the size of the trade?
\item How do \omterm{def:m2:how-bonds-trade:rfq}{axes} change the choice?
\item How would \omterm{def:m2:how-bonds-trade:a2a}{all-to-all trading} change it?
\end{enumerate}

\medskip\noindent\textbf{Part III --- Alternatives.}
\begin{enumerate}[resume]
\item When would the manager include the bond in a \omterm{def:m2:how-bonds-trade:a2a}{portfolio trade} instead?
\item What does TRACE show of the sale, and when?
\item How does the \omterm{def:m2:how-bonds-trade:trace}{dissemination cap} protect the winning dealer?
\item Could the manager sell ETF shares instead of the bond?
\item What does a \omterm{def:m2:how-bonds-trade:composite}{composite price} tell the manager before it asks?
\end{enumerate}

\medskip\noindent\textbf{Part IV --- Judgement.}
\begin{enumerate}[resume]
\item Why do dealers not always bid their best to every client?
\item How would you measure leakage in your own trading records?
\item Why might the best number differ for a high-yield bond?
\item State the \emph{named result}: the number of dealers that minimises the expected
cost, and that cost.
\item In one sentence: why is asking more dealers not always better?
\end{enumerate}
\end{problem}

\section{Interview questions}

\begin{interviewq}[$\star$ \iqroles{trader, developer}]\label{iq:m2:how-bonds-trade:1}
How is a corporate bond bought and sold, compared with a stock?
\end{interviewq}

\begin{interviewq}[$\star$ \iqroles{trader, researcher}]\label{iq:m2:how-bonds-trade:2}
What is TRACE, and what does it not show?
\end{interviewq}

\begin{interviewq}[$\star\star$ \iqroles{researcher}]\label{iq:m2:how-bonds-trade:3}
How many dealers should a client ask in an RFQ? Build a simple model.
\end{interviewq}

\begin{interviewq}[$\star\star$ \iqroles{trader, bank}]\label{iq:m2:how-bonds-trade:4}
How does a dealer price a \omterm{def:m2:how-bonds-trade:a2a}{portfolio trade} of 400 bonds?
\end{interviewq}

\begin{interviewq}[$\star\star$ \iqroles{researcher, trader}]\label{iq:m2:how-bonds-trade:5}
Why can a bond ETF trade at a discount to its net asset value, and which price is
right?
\end{interviewq}

\begin{interviewq}[$\star\star\star$ \iqroles{developer}]\label{iq:m2:how-bonds-trade:6}
Design a system that answers 20\,000 bond RFQs a day with prices within 100
milliseconds.
\end{interviewq}
